\section{Leaderboard Success 为什么不等于 Commonsense}

Maieutic 的提升仍发生在 benchmark 内部，因此讲者下一步重新审视 benchmark 本身。如果数据集存在固定 artifact，模型可能学会“如何拿分”而没有学会目标能力；这就是 dataset solving 与 task solving 的差别。

\subsection{高分与脆弱性可以同时存在}

slide 将 object recognition、question answering 和 captioning 的高分结果，与 adversarial 或 out-of-domain 失败放在一起。横向比较不是为了否定进展，而是指出评价分布通常比真实使用窄。底部结论“只解决 dataset，而未解决 underlying task”要求研究者报告 perturbation、domain shift 与 failure category，而不是只报 leaderboard rank。

\lecturefigure{slide-09-brittle-leaderboards-dataset-task.jpg}{Superhuman leaderboard performance 与 adversarial brittleness 并存}{官方录像恢复幻灯片 V009；Stanford CS25 Lecture 16}

\teachervoice{讲者的“choose your battles”在这里落地：如果 benchmark 无法区分 surface pattern 与 conceptual understanding，继续堆分只是在错误战场获胜。新的 task 设计本身就是算法研究的一部分。}

\subsection{Human 与 AI 之间缺的不是更多句子}

桥接图把人类日常生活、语言与社会经验放在左侧，把多彩但不稳定的机器输出放在右侧。中间的 gap 指向 practical knowledge：知道冰箱门通常应关闭、知道一句话在特定关系中为何冒犯、知道动作可能造成什么物理后果。模型需要的不只是更多 text tokens，而是可组合、可检验的情境知识。

\lecturefigure{slide-10-commonsense-gap.jpg}{Commonsense gap：从表面语言到日常世界理解}{官方录像恢复幻灯片 V010；Stanford CS25 Lecture 16}

\subsection{本讲采用的 Operational Definition}

定义 slide 把 commonsense 限定为 practical knowledge and reasoning，关注 everyday situations and events，并且通常由大多数人共享。三个限定都重要：它不是百科知识全集，不只是一组 fact，也不是绝对跨文化一致。读图时要把“commonly shared”理解为经验性范围，而非道德真理或 universal law。

\lecturefigure{slide-11-definition-of-commonsense.jpg}{本讲的常识定义：日常、实用、广泛共享的知识与推理}{官方录像恢复幻灯片 V011；Stanford CS25 Lecture 16}

下一张 slide 进一步区分人类和机器的用途。对人类，常识支持共处与安全；对 AI，常识帮助理解用户需求、行动后果与隐含条件。冰箱门例子表明同一句“保持门打开”在不同情境下可能合理或危险，因此常识不是固定 phrase lookup，而是 context-conditioned inference。

\lecturefigure{slide-12-human-ai-commonsense-importance.jpg}{常识为何同时对人类协调和 AI 行动理解重要}{官方录像恢复幻灯片 V012；Stanford CS25 Lecture 16}

\begin{importantbox}{Dataset、Task 与 Capability 三层}
Dataset 是有限样本；task 是希望系统完成的映射；capability 是跨分布复用的稳定能力。测试集高分最多直接证明 dataset-level performance。要主张 commonsense capability，还需要多样改写、组合变化、反事实、跨域和 failure analysis。
\end{importantbox}

\subsection{本章小结}

常识研究的第一项难题是定义评价对象。排行榜并非无用，但必须与 adversarial 和 out-of-domain evidence 配对。后续 ATOMIC、COMET 与 Delphi 的意义，正是把“模型也许隐含知道”转化成更明确的数据、关系和任务接口。

